

We saw, both in our case study on the Atari 100k benchmark and with our analysis of other widely-used RL benchmarks, that statistical issues can have a sizeable influence on reported results, in particular when point estimates are used or evaluation protocols are not kept constant within comparisons. 
Despite earlier calls for more experimental rigor in deep RL~\citep{henderson2018deep, colas2018many, colas2019hitchhiker, jordan2020evaluating, chan2019measuring,recht-2018}~(discussed in Appendix~\ref{related_work}), our analysis shows that the field has not yet found sure footing in this regards. 

In part, this is because the issue of reproducibility is a complex one; where our work is concerned with our confidence about and interpretation of reported results (what \citet{goodman2016does} calls \emph{results reproducibility}), others~\citep{pineau2020improving} have highlighted that there might be missing information about the experiments themselves~(\emph{methods reproducibility}). We remark that the problem is not solved by fixing random seeds, as has sometimes been proposed~\citep{nagarajan2018deterministic, kolesnikov2019catalyst}, since it does not really address the question of whether an algorithm would perform well under similar conditions but with different seeds. Furthermore, fixed seeds might benefit certain algorithms more than others. Nor can the problem be solved by the use of dichotomous statistical significance tests, as discussed in \Secref{sec:formalism}.

One way to minimize the risks associated with statistical effects is to report results in a more complete fashion, paying close attention to bias and uncertainty within these estimates. To this end, our recommendations are summarized in \tabref{tab:recommendations}.
To further support RL researchers in this endeavour, we released an easy-to-use Python library, \href{https://github.com/google-research/rliable}{\ttfamily rliable} along with a \href{https://bit.ly/drl\_precipice\_colab}{Colab notebook} for implementing our recommendations, as well as all the individual \href{https://console.cloud.google.com/storage/browser/rl-benchmark-data/}{runs} used in our experiments\footnote{Colab: \href{https://bit.ly/statistical\_precipice\_colab}{\ttfamily  bit.ly/statistical\_precipice\_colab}. Individual runs: \href{https://console.cloud.google.com/storage/browser/rl-benchmark-data/}{\ttfamily gs://rl-benchmark-data}.}. Again, we emphasize the importance of published papers providing results for all runs to allow for future statistical analyses.

A barrier to adoption of evaluation protocols proposed in this work, and more generally, rigorous evaluation, is whether there are clear incentives for researchers to do so, as more rigor generally entails more nuanced and tempered claims. Arguably, doing good and reproducible science is one such incentive. We hope that our findings about erroneous conclusions in published papers would encourage researchers to avoid fooling themselves, even if that requires tempered claims. That said, a more pragmatic incentive would be if conferences and reviewers required more rigorous evaluation for publication, \eg NeurIPS 2021 checklist asks whether error bars are reported. Moving towards reliable evaluation is an ongoing process and we believe that this paper would greatly benefit it.


Given the substantial influence of statistical considerations in experiments involving 40-year old Atari 2600 video games and low-DOF robotic simulations, we argue that it is unlikely that an increase in available computation will resolve the problem for the future generation of RL benchmarks. Instead, just as a well-prepared rock-climber can skirt the edge of the steepest precipices, it seems likely that ongoing progress in reinforcement learning will require greater experimental discipline.


\section*{Societal Impacts}
\vspace{-0.1cm}
This paper calls for statistical sophistication in deep RL research by accounting for statistical uncertainty in reported results. However, statistical sophistication can introduce new forms of statistical abuses and monitoring the literature for such abuses should be an ongoing priority for the research community.
Moving towards reliable evaluation and reproducible research is an ongoing process and this paper only partly addresses it by providing tools for more reliable evaluation. That said, while accounting for uncertainty in results is not a panacea, it provides a strong foundation for trustworthy results on which the community can build upon, with increased confidence. In terms of broader societal impact of this work, we do not see any foreseeable strongly negative impacts. However, this paper could positively impact society by constituting a step forwards in rigorous few-run evaluation regime, which reduces computational burden on researchers and is ``greener'' than evaluating a large number of runs. 

\vspace{-0.1cm}
\section*{Acknowledgments} 
\vspace{-0.1cm}
We thank Xavier Bouthillier, Dumitru Erhan, Marlos C. Machado, David Ha, Fabio Viola, Fernando Diaz,  Stephanie Chan, Jacob Buckman, Danijar Hafner and anonymous NeurIPS' reviewers for providing valuable feedback for an earlier draft of this work. We also acknowledge  Matteo Hessel, David Silver, Tom Schaul, Csaba Szepesvári, Hado van Hasselt, Rosanne Liu, Simon Kornblith, Aviral Kumar, George Tucker, Kevin Murphy, Ankit Anand, Aravind Srinivas, Matthew Botvinick, Clare Lyle, Kimin Lee, Misha Laskin, Ankesh Anand, Joelle Pineau and Braham Synder for helpful discussions. We also thank all the authors who provided individual runs for their corresponding publications. We are also grateful for general support from Google Research teams in Montréal and elsewhere.